% *********************************************************************
% © 2016–2025 Jeremy Sylvestre
%
% Permission is granted to copy, distribute and/or modify this document
% under the terms of the GNU Free Documentation License, Version 1.3 or
% any later version published by the Free Software Foundation; with no
% Invariant Sections, no Front-Cover Texts, and no Back-Cover Texts. A
% copy of the license is included in the appendix entitled “GNU Free
% Documentation License” that appears in the output document of this
% PreTeXt source code. All trademarks™ are the registered® marks of
% their respective owners.
%
% *********************************************************************
\tdplotsetmaincoords{95}{30}
\begin{tikzpicture}[
	point/.style={circle,draw,very thin,fill,inner sep=0pt,minimum size=4pt},
	vector/.style={-latex},
	tdplot_main_coords
]

	% plane through the origin with normal vector (1,-1,2)
	% orthog basis used to draw sides: (1,1,0) and (1,-1,-1)
	% note that the origin is not at (0,0,0): point x0 acts as an "origin" on the plane

	% origin
	\node[point] at (-1,-1,-2.5) (o) [label=below left:{$\zerovec$}] {};
	% plane corners
	% sides of plane are vectors (6,6,0) and (2.5,-2.5,-2.5)
	\coordinate (P1) at (-1.5,-0.5,0.5);
	\coordinate (P2) at (4,5,0.5);
	\coordinate (P3) at (6.5,2.5,-2);
	\coordinate (P4) at (1,-3,-2);

	% fill/border for plane
	\filldraw[opacity=0.25,gray] (P1) to (P2) to (P3) to (P4) to cycle;
	\draw[dotted]  (P1) to (P2) to (P3) to (P4) to cycle;

	\begin{scope}[opacity=0.3]

		% grid lines parallel to p1, spaced by p2
		\draw (-0.435-1,-0.58,-0.0725+0.5) to (-0.075-1,-0.1,-0.0125+0.5);
		\draw (-0.8625,-1.15,-0.14375) to (2.925,3.9,0.4875); % middle
		\draw (-1.275+1,-1.7,-0.2125-0.5) to (3.435+1,4.58,0.5725-0.5);
		\draw (-1.725+2,-2.3,-0.2875-1) to (3+2,4,0.5-1);
		\draw (-2.1375+3,-2.85,-0.35625-1.5) to (2.58+3,3.44,0.43-1.5);
		\draw (0+4,0,0-2) to (2.145+4,2.86,0.3575-2);

		% grid lines parallel to p2, spaced by p1
		\draw (1.5-1.5,0-2,-0.75-0.25) to (3.5-1.5,0-2,-1.75-0.25);
		\draw (-0.25-0.75,0-1,0.125-0.125) to (3.75-0.75,0-1,-1.875-0.125);
		\draw (-1,0,0.5) to (4,0,-2); % middle
		\draw (-0.75+0.75,0+1,0.375+0.125) to (4.25+0.75,0+1,-2.125+0.125);
		\draw (-0.5+1.5,0+2,0.25+0.25) to (4.5+1.5,0+2,-2.25+0.25);
		\draw (-0.25+2.25,0+3,0.125+0.375) to (3.75+2.25,0+3,-1.875+0.375);
		\draw (0+3,0+4,0+0.5) to (2+3,0+4,-1+0.5);

	\end{scope}

	% plane base point
	\node[point] at (0,0,0) (x0) [label=left:$X_0$] {};
	% base point vector transition point = 60% * (o) (viewed in opposite direction from drawn vector)
	\coordinate (x0_trans) at (-0.6,-0.6,-1.5);
	% plane variable point
	\node[point] at (5.25,3,-1.125) (x) [label=below:$X$] {};
	% vector (o) to (x) is (6.25,4,1.375)
	% variable point vector transition point = (o) + 47% * (x - o)
	\coordinate (x_trans) at (1.938,0.88,-1.854);
	% plane basis vector terminal points
	\coordinate (p1) at (0.75,1,0.125);
	\coordinate (p2) at (1,0,-0.5);
	% smul of plane basis vector terminal points
	\coordinate (sp1) at (2.25,3,0.375);
	\coordinate (tp2) at (3,0,-1.5);

	% vector for base point
	\draw[thick] (o) to node[above left] {$\uvec{x}_0$} (x0_trans);
	\draw[vector,thick,dashed] (x0_trans) to (x0);

	% extended vectors
	\draw[vector] (x0) to node[above,at end] {$s \uvec{p}_1$} (sp1);
	\draw[vector] (x0) to node[below,near end] {$t \uvec{p}_2$} (tp2);

	% vector for "variable" point
	\draw[thick] (o) to node[below, near end, xshift=1cm] {$\uvec{x} = \uvec{x}_0 + s \uvec{p}_1 + t \uvec{p}_2$} (x_trans);
	\draw[vector,dashed,thick] (x_trans) to (x);

	% plane vectors
	\draw[vector,thick] (x0) to node[above left,near end] {$\uvec{p}_1$} (p1);
	\draw[vector,thick] (x0) to node[below,near end,yshift=-0.1cm] {$\uvec{p}_2$} (p2);
	\draw[vector,thick] (x0) to node[sloped,above] {$s \uvec{p}_1 + t \uvec{p}_2$} (x);

\end{tikzpicture}
